\chapter{Sub-100 digits: a reduction to two hard problems}\label{ch:barriers}

Chapter~\ref{ch:frontier} priced the threshold game at $Q=10^{5}$ and found the
$100$-digit rung seven orders of magnitude beyond plausible compute. This
chapter asks whether that wall is real. The answer is that it is real
\emph{within the paradigm}, and that stepping outside the paradigm leads
directly to two well-known hard problems---which is a more useful outcome than
either a proof or a construction, because it says exactly what would have to be
broken.

It is worth being clear at the outset that sub-$100$-digit deserts \emph{exist}
in abundance. With no cover at all, the failure exponent at $B=10^{100}$ is
$E_{\text{random}}\approx110$, so the density of qualifying integers is
$e^{-110}$ and the expected count below $10^{100}$ is about
$10^{100}\cdot e^{-110}\approx10^{52}$. Nothing here is an existence
obstruction. The entire difficulty is constructive.

\section{Two regimes}

Write $A=\log M$ for the cover mass in nats and $\log B$ for the digit budget
($\log B\approx230$ at $100$ digits). Covers split into two regimes according to
whether their modulus fits under the budget.

\paragraph{Regime I: $A\le\log B$.} The CRT progression itself supplies
$e^{\log B-A}$ candidates, and the search costs $e^{E}$. Optimising at
$Q=10^{5}$, $B=10^{100}$ gives a best cover of $44$ moduli with
$M\approx10^{79}$, $|U|=1060$, $\mathrm{boost}=9.6$, hence
\[
E\;=\;44.1,\qquad e^{E}\approx1.4\times10^{19}\ \text{candidates}.
\]
At the record engine's measured throughput of about $5\times10^{5}$ candidates
per second this is roughly $9\times10^{5}$ GPU-years.

\begin{remark}[A units error, corrected]\label{rem:units}
An earlier version of this calculation reported $5\times10^{4}$ GPU-years, an
error of a factor of $18$: it divided the candidate count by the \emph{test}
rate ($8.5\times10^{6}$ Fermat tests/s) rather than the \emph{candidate} rate,
forgetting that a candidate costs about $17$ tests. The error was caught by
cross-checking against an independent study of the same frontier, which measured
a realisable $E=43.66$ at $100$ digits---within $0.4$ nats of the figure above,
and irreconcilable with the faster cost. We record the correction here because
the discipline that caught it (always cross-check a headline cost against an
independently derived one) is worth more than the number.
\end{remark}

\paragraph{Regime II: $A>\log B$ (oversized covers).} If the modulus is allowed
to exceed the budget, coverage becomes excellent. All $9\,592$ primes below
$10^{5}$ can be covered by $355$ moduli of total mass $1\,019$ digits; allowing
each modulus to use $L=16$ residue classes, $420$ moduli leave only $|U|=165$,
giving
\[
E\;=\;10.2,\qquad e^{E}\approx27\,000\ \text{candidates}.
\]
Twenty-seven thousand candidates is nothing. If we could produce such a design
\emph{together with} a representative below $10^{100}$, sub-$100$ would fall
this afternoon.

Solutions exist: the design space (which moduli, which classes) carries about
$2\,698$ nats of entropy against $\log(M/B)=2\,630$ nats of ``distance to
squeeze'', leaving a surplus of about $+57$ nats and hence roughly $e^{57}$
valid (design, $\N$) pairs below $10^{100}$. But the \emph{density} of designs
whose representative lands below $B$ is $e^{-2630}$. Enumeration is out by some
$1\,100$ orders of magnitude.

\begin{observation}
Regime II is where sub-$100$ becomes easy \emph{if and only if} small
representatives can be found directly. That is the whole problem.
\end{observation}

\section{The small-representative problem is modular subset-sum}

CRT reconstruction is additive. Writing the standard idempotent basis,
\[
\N\;\equiv\;\sum_{r\in S}a_r\cdot\frac{M}{r}\cdot\Bigl(\bigl(\tfrac{M}{r}\bigr)^{-1}\bmod r\Bigr)
\pmod M ,
\]
so ``find a design whose representative is below $B$'' is exactly:
\begin{quote}\itshape
minimise $\bigl|\sum_r c_r \bmod M\bigr|$ over choices $c_r\in C_r$ with
$|C_r|=L$, targeting the window $[0,B)$.
\end{quote}
This is an inhomogeneous \emph{modular subset-sum problem with per-position
choices}. Its density---the ratio of available design entropy to the window
squeeze---is
\[
d\;=\;\frac{2\,698}{2\,630}\;\approx\;1.026 .
\]

That is the worst possible value. Low-density instances ($d<0.94$) fall to
lattice reduction~\cite{LagariasOdlyzko,Coster}; very high density instances
fall to birthday and $k$-tree methods~\cite{Wagner}. Density $\approx1$ is
precisely the regime where neither works and where the presumed hardness of
subset-sum is concentrated. Concretely, our solution set has size about
$2^{82}$ inside a design space of about $2^{3893}$; meet-in-the-middle gives
$2^{1897}$ and Wagner's $k$-tree with our list sizes still needs $2^{938}$.
Both are astronomically short.

\section{The decoding shortcut, and why it fails}

There is one route that would bypass the subset-sum problem entirely. Finding a
small integer consistent with many congruences is a decoding problem for the
Chinese Remainder code, and CRT codes have list-decoding
algorithms~\cite{GSS,Goldreich}. If the decoder's radius were large enough, we
could simply ask it for a small $\N$ agreeing with a chosen cover.

The relevant statement is Guruswami--Sahai--Sudan~\cite{GSS}, Theorem 3: for a
CRT code with moduli $p_i$, message bound $K$, and any non-negative integers
$\ell$ and $z_i$, one finds in time $\mathrm{poly}(n,\log N,\ell,\sum z_i)$ all
codewords $m$ with
\[
\sum_i \alpha_i z_i\log p_i\;>\;\log(\ell+1)+\frac{\ell}{2}\log\frac{K}{2}
+\frac{1}{\ell+1}\sum_i\binom{z_i+1}{2}\log p_i ,
\]
where $\alpha_i=1$ exactly when $m\equiv r_i\pmod{p_i}$. Optimising $z$ and
$\ell$ under uniform weights, with $L$ allowed classes per position, gives an
agreement-mass radius
\begin{equation}\label{eq:gss}
A\;>\;\sqrt{L\cdot\log B\cdot P},\qquad P=\sum_{\text{pool}}\log p_i .
\end{equation}
We verified~\eqref{eq:gss} numerically against the exact finite-$\ell$ optimum:
the ratio is $1.001$--$1.004$ across pools from $5\times10^{3}$ to $10^{5}$ and
$L\in\{1,2\}$.

Now the trap closes. The radius scales with the mass $P$ of the \emph{whole
pool}, which must be far larger than the cover itself in order to supply the
subset entropy that makes small representatives exist at all. Numerically, at
pool bound $60\,000$ with $L=16$: the pool mass is $59\,816$ nats, the decoding
radius is $14\,845$ nats, and our cover mass is $2\,860$ nats. The decoder needs
agreement $5.2\times$ larger than anything we can supply. Shrinking the pool
lowers the radius but destroys the entropy that makes solutions exist.

\begin{theorem}[Barrier; heuristic, first-moment]\label{thm:barrier}
For CRT-cover constructions of even $\N<B$ with $\g(\N)>Q$, no parameter choice
simultaneously satisfies
\begin{enumerate}
\item[(a)] $A<\log B+\log\binom{n}{s}+s\log L$ \quad (representatives exist), and
\item[(b)] $A>\sqrt{L\cdot\log B\cdot P}$ \quad (GSS-decodable).
\end{enumerate}
Hence sub-$100$ via oversized covers requires either list-recovery of CRT codes
\emph{beyond} the Johnson-type radius---an open problem in coding theory---or a
subset-sum algorithm at density $\approx1$.
\end{theorem}

\begin{table}[h]
\centering\small
\caption{The barrier, evaluated. ``Exists'' is the largest cover mass for which
representatives below $B$ still exist; ``decodes'' is the GSS radius. The ratio
never approaches $1$.}\label{tab:barrier}
\begin{tabular}{rrrrrr}
\toprule
pool bound & $L$ & pool mass & $A$ exists & $A$ decodes & ratio\\
\midrule
$1\,000$ & $1$ & $956$ & $339$ & $469$ & $0.72$\\
$60\,000$ & $1$ & $59\,816$ & $570$ & $3\,711$ & $0.15$\\
$60\,000$ & $16$ & $59\,816$ & $1\,305$ & $14\,845$ & $0.09$\\
$10^{6}$ & $16$ & $998\,483$ & $1\,697$ & $60\,651$ & $0.03$\\
\bottomrule
\end{tabular}
\end{table}

\subsection{A retracted claim, and what it teaches}\label{sec:retraction}

An earlier version of this analysis reported an \emph{open window}: a range of
pool sizes where, apparently, representatives existed and the decoder still
worked. That claim was wrong and is withdrawn. The error is instructive enough
to record.

The two conditions of Theorem~\ref{thm:barrier} both depend on an agreement
size $s$---how many positions the sought codeword matches. The retracted
analysis computed the decoding radius at one value of $s$ and the existence
bound at a \emph{different} value (the one maximising existence, near $2n/3$),
and compared the two numbers. Requiring both conditions at the \emph{same} $s$
closes the window at every pool size. The lesson generalises well beyond this
problem: \emph{a feasibility window assembled from two separately optimised
bounds is worthless; couple the parameters first, then optimise once.}

We also had the barrier re-derived independently, by a route that shares no
code with ours---a Legendre-transform treatment of the same optimisation---which
reproduces the $14\,845$-nat radius and finds the gap closed at every pool size,
with the best case still $28\%$ short. Independent confirmation matters
especially when the first version of a result was wrong.

\section{Does lattice reduction beat the radius in practice?}

Decoding radii are worst-case statements. Our instances are highly structured
and, as computed above, admit roughly $e^{57}$ solutions---so it is entirely
reasonable to hope that lattice reduction succeeds far beyond the proved radius,
as it often does elsewhere. We tested this directly rather than assuming either
way.

We built a small CRT-decoding testbed: plant a known small solution in a
random instance at a controlled agreement mass, then run
LLL~\cite{LLL} on the standard Howgrave-Graham-style lattice and record whether
the planted solution (or any other) is recovered. Sweeping the agreement mass
across the Johnson-type radius $A_J$ gives a sharp threshold:

\begin{itemize}
\item at $\ell=30$, recovery succeeds above $1.028\,A_J$ and fails below;
\item at $\ell=45$ the threshold tightens to about $1.02\,A_J$, converging to
the theoretical radius from above;
\item below the threshold the failure is \emph{total}: across $390$ decodes we
recovered zero non-planted roots, and no partial or near-solutions.
\end{itemize}

The last point is the one that matters. Solution abundance buys nothing: even
with $10^{7}$ alternative solutions available in the instance, the lattice
returns nothing at all below the radius. The Johnson-type radius behaves here
like a genuine algorithmic boundary rather than a proof artifact, which is what
makes Theorem~\ref{thm:barrier} a barrier and not merely a gap in our knowledge.

\section{Where this leaves sub-100}

\begin{enumerate}
\item \textbf{Regime I is the only currently viable route}, and it is out of
reach: $E=44.1$, about $9\times10^{5}$ GPU-years, seven orders of magnitude
beyond plausible compute. The fundable frontier is $130$--$140$ digits
(Table~\ref{tab:menu}).

\item \textbf{Two well-posed open problems now sit under the record}, either of
which would collapse it: (i) list-recovery of CRT codes beyond the Johnson-type
radius for structured, solution-abundant instances; (ii) modular subset-sum with
per-position choices at density $\approx1$ where $2^{82}$ solutions exist. Both
are of independent interest, which is the right way for a chapter like this to
end.

\item \textbf{The one place the barrier does not apply} is a construction that
certifies compositeness by something other than a congruence divisor. Every
argument in this chapter accounts entropy in units of $\log r$ per excluded
offset; a mechanism that escapes that accounting escapes the barrier. Whether
any such mechanism exists is the subject of Chapter~\ref{ch:closure}, and the
answer, within a precisely defined language, is no.
\end{enumerate}
